\chapter{Analysis of the \texttt{linkfigs.m} MATLAB Script}
\date{\today}


\maketitle

\section{Overview}

\subsection*{What is the main purpose of this file?}
The primary purpose of the \texttt{linkfigs.m} script is to generate sets of coordinates and their corresponding topological connections for various geometric shapes, numbers, and letters. It acts as a data generator for creating custom, parameterized node-link structures.

\subsection*{What type of analysis or calculation does it perform?}
The script performs geometric and topological definitions. It calculates explicit coordinates in the complex plane (using complex numbers where the real part is the x-axis and the imaginary part is the y-axis) and defines a mapping array that determines how these coordinates are connected sequentially.

\subsection*{What is the mathematical/theoretical context?}
The theoretical context lies within the realm of \textbf{discrete mathematics, directed graphs, and endomorphisms on finite sets}. 
The array $g \in \mathbb{C}^n$ represents an embedding of vertices in the complex plane. The array $\phi$ (phi) represents an endomorphism $f: V \to V$ on the set of vertices, where the $i$-th element of $\phi$ points to the index of the target vertex. This maps every vertex to exactly one other vertex, creating directed paths, cycles (as in the case of 'A' or '1'), or fixed points (as in the case of 'M' and 'V'). Furthermore, the '*' and 'circles' cases explore the complex exponential map of a parameterized grid, essentially defining discrete approximations of Riemann surfaces or parameterized manifolds.

\section{Step-by-Step Analysis}

\subsection{1. Function Definition and Initialization}
\textbf{Explanation:} This section defines the function signature. The function takes two arguments: \texttt{Fig}, a string determining the shape to generate, and \texttt{w}, a complex number that acts as a scaling, translation, or parameterization factor depending on the chosen shape. It returns \texttt{g} (the complex coordinates) and \texttt{phi} (the connectivity mapping). A \texttt{switch} statement is initiated to handle the different \texttt{Fig} inputs.

\begin{lstlisting}[language=Matlab]
function [g,phi]=linkfigs(Fig,w)
% 14Oct2024 changes output to [phi,g]
  % [phi,g]=linkfigures(Fig,w)
  %
  % Returns a link structure for figures in the database, currently letter A and Number 1
  % added a star '*'

  switch Fig
\end{lstlisting}

\subsection{2. Case 'A' and '1' (Polygonal Shapes)}
\textbf{Explanation:} These cases define the coordinates for the letter 'A' and the number '1'. The \texttt{phi} array \texttt{[2:8,1]} creates a closed loop (a cycle of length 8), where node 1 connects to 2, 2 to 3, ..., and 8 connects back to 1. The \texttt{g} array calculates the precise vertices in the complex plane. The parameter \texttt{w} is used to manipulate the geometry (e.g., scaling or shifting certain vertices).

\begin{lstlisting}[language=Matlab]
    case 'A'
      phi=[2:8,1];
      g=[-1; 0; w; w-1; w/sqrt(2)-1; -conj(w)/sqrt(2); -conj(w); -conj(w)-1];
    case '1'
    phi=[2:8,1]; b=imag(w);
    g=[-1; 0; w; b*1i; pi*b*1i; pi*b*1i + w/pi; -1+pi*b*1i-conj(w)/pi; -1+pi*b*1i ];
\end{lstlisting}

\subsection{3. Case '*' and 'circles' (Exponential Grids)}
\textbf{Explanation:} These cases build a 2D grid of nodes using \texttt{ndgrid}. The grid dimensions $n$ and $m$ are extracted from the real and imaginary parts of \texttt{w}. A logarithmic mapping is applied to the real axis and a linear mapping to the imaginary axis (from 0 to $2\pi$). The coordinates $g$ are generated by passing these values through the complex exponential function $g = \exp(t)$, wrapping the grid onto a cylindrical or circular topology. The mapping \texttt{phi} is manipulated to connect adjacent columns, with 'circles' wrapping the mapping continuously around the structure.

\begin{lstlisting}[language=Matlab]
  case '*'
      n=ceil(real(w)); m=ceil(imag(w));
      phi=reshape((1:n*m),[n m]);
      phi(1,:)=[]; phi(end+1,:)=phi(end,:); phi=phi(:);
      [a,b]=ndgrid(log((1:n)/n),(1:m)/m*2*pi);
      t=complex(a,b);
      g=exp(t(:));
    case 'circles'
      n=ceil(real(w)); m=ceil(imag(w))+1;
      phi=reshape((1:n*m),[n m]);
      phi=reshape(phi(:,[2:end 1]),[m*n 1]);
      [a,b]=ndgrid(log((1:n)/n),(1:m)/m*2*pi);
      t=complex(a,b);
      g=exp(t(:));
\end{lstlisting}

\subsection{4. Case 'M' and 'V' (Letters with Fixed Points)}
\textbf{Explanation:} Here, coordinates for the letters 'M' and 'V' are defined manually using complex numbers. Notice the \texttt{phi} array: \texttt{[2 3 4 5 5]}. Node 1 connects to 2, 2 to 3, 3 to 4, 4 to 5, and 5 connects to itself. Node 5 is a \textit{fixed point} under this endomorphism, meaning the drawing path terminates or rests there rather than looping indefinitely.

\begin{lstlisting}[language=Matlab]
      case 'M'
      phi=[2 3 4 5 5];
      g=[0+0i, 0+1i, 0.5+0.5i, 1+1i, 1+0i];
    case 'V'
      phi=[2 3 4 5 5];
      g=[0+1i, 0+0.5i, 0.5+0i, 1+0.5i, 1+1i];
\end{lstlisting}

\subsection{5. Case 'line' (Linear Interpolation)}
\textbf{Explanation:} This case generates a straight line of $n$ points spaced evenly from 0 to 1 along the real axis using \texttt{linspace}. The number of points $n$ and a connection parameter $m$ are derived from \texttt{w}. The \texttt{phi} array connects nodes sequentially, but the last node connects to node $m$, creating a "lasso" or a line that loops back into a cycle at a specific internal node.

\begin{lstlisting}[language=Matlab]
    case 'line'  % w=5+5i
      n=abs(ceil(real(w))); m=abs(ceil(imag(w)));
      if m>n, m=n; end
      if m<1, m=1; end
      phi=[2:n, m];
      g=linspace(0,1,n);
\end{lstlisting}

\subsection{6. Default Case (Error Handling)}
\textbf{Explanation:} If the user inputs an unrecognized string for \texttt{Fig}, the script alerts the user with a console message detailing the valid options and returns empty arrays for both \texttt{g} and \texttt{phi} to prevent downstream errors.

\begin{lstlisting}[language=Matlab]
    otherwise
      disp([Fig ' is not a valid input string'])
      disp('The current possibilities are: A, 1, *, circles')
      phi=[]; g=[];
  end
\end{lstlisting}

\section{Expected Outputs and Visuals}

While the \texttt{linkfigs.m} script does not invoke plotting functions (like \texttt{plot} or \texttt{gplot}) directly, it is clearly intended to feed into a visualizer. If one were to plot the outputs (e.g., using a custom loop to draw lines from \texttt{g(i)} to \texttt{g(phi(i))}), the expected visuals would be:

\begin{itemize}
    \item \textbf{Axes:} The horizontal axis represents the Real part of the coordinates ($\text{Re}(g)$), and the vertical axis represents the Imaginary part ($\text{Im}(g)$).
    \item \textbf{Polygons ('A', '1', 'M', 'V'):} The plot would yield 2D line drawings resembling the respective characters. For 'A' and '1', it would trace a closed polygonal loop outlining the character. For 'M' and 'V', it would draw an open path from the first node to the final fixed node.
    \item \textbf{Exponential Maps ('*', 'circles'):} Plotting these would result in visually rich, concentric circular or spiral grid webs. Because $g = \exp(a + bi) = e^a(\cos(b) + i\sin(b))$, the points form concentric rings (since $a = \log(r)$ bounds the radius) distributed angularly from $0$ to $2\pi$. The mapping \texttt{phi} would draw a spiderweb-like directed graph over these nodes.
    \item \textbf{Line ('line'):} A straight horizontal line segmented into $n$ points along the x-axis from $x=0$ to $x=1$. The final node would have a directed edge jumping back to the $m$-th node.
\end{itemize}

\vspace{1em}
\noindent \textit{Would you like me to write a corresponding MATLAB wrapper script that takes the outputs of \texttt{linkfigs} and plots these figures visually?}
